\documentclass[11pt]{article}
\usepackage[margin=1in]{geometry}
\usepackage{amsmath,amssymb,amsthm}
\usepackage{booktabs}
\usepackage{graphicx}
\usepackage{array}
\usepackage{siunitx}
\usepackage[authoryear,round]{natbib}
\usepackage[hidelinks]{hyperref}
\usepackage{authblk}
\usepackage{caption}
\captionsetup{font=small,labelfont=bf}
\newtheorem{proposition}{Proposition}
\newcommand{\R}{\mathbb{R}}
\newcommand{\E}{\mathbb{E}}
\newcommand{\Prob}{\mathbb{P}}
\newcommand{\VaR}{\mathrm{VaR}}
\newcommand{\ES}{\mathrm{ES}}
\newcommand{\TER}{\mathrm{TER}}
\newcommand{\CTER}{\mathrm{CTER}}
\newcommand{\ind}{\mathbf{1}}
\DeclareMathOperator*{\argmax}{arg\,max}

\title{\bfseries Conditional Tail Episode Risk:\\ A Path-Dependent Framework for Extreme-Loss Episodes\\ Beyond Value-at-Risk and Expected Shortfall}
\author{Niraj Neupane, CA (ICAI)}
\affil{\normalsize Quantitative Researcher, Korvane \quad$\cdot$\quad Financial Economist, Calderyn Institute\\
\normalsize ORCID: 0009-0003-7026-7026 \quad$\cdot$\quad \texttt{nirajneupane17@gmail.com}}
\date{Working Paper --- This version: September 21, 2026}

\begin{document}
\maketitle

\begin{abstract}
\noindent Value-at-Risk (VaR) and Expected Shortfall (ES) are central tools for financial risk
measurement, but both primarily characterize the \emph{marginal} distribution of losses. Consequently
they can assign similar risk to two loss paths that contain the same individual tail observations but
differ in the temporal organization of those observations. This paper develops \textbf{Tail Episode Risk
(TER)}, a path-dependent tail-risk functional that quantifies the severity of the most extreme
\emph{contiguous} cumulative-loss episode within a forecast horizon, and \textbf{Conditional Tail Episode
Risk (CTER)}, which decomposes future episode risk into (i) the conditional probability that an
extreme-loss episode occurs and (ii) the conditional upper-tail severity of that episode. The framework is
related to aggregate excess severity in extreme-value theory, clustered extremes, Conditional Expected
Drawdown, and dynamic VaR/ES forecasting, but targets a distinct state variable: cumulative loss above a
tail threshold within a contiguous episode.

We establish basic properties of TER---monotonicity, positive homogeneity, translation invariance
under a translation-equivariant threshold rule (in the sense of a severity/deviation functional rather than
a monetary risk measure), and sensitivity to temporal clustering---and give constructed
counterexamples in which two paths share identical marginal VaR and ES but differ in TER. Monte Carlo
experiments confirm that TER increases sharply with serial dependence while marginal VaR and ES remain
essentially flat, and adversarial stress tests show that TER responds to sustained episodes rather than to
isolated tail observations.

Empirically, we evaluate the framework on daily data for six equity, bond, foreign-exchange, and
cryptocurrency markets over 2000--2026, with asset-specific out-of-sample windows (from 2005 for equities
and FX, 2008 for Treasuries, and 2018 for Bitcoin). The evidence is \emph{cautionary}. Under
a strictly consistent Fissler--Ziegel scoring rule, GARCH-$t$ is the strongest conventional benchmark and
an EWMA specification a strong simple one. Episode-\emph{occurrence} discrimination is weak and decays into
the tail (mean AUC $0.56$, $0.54$, $0.47$ at the $95\%$, $97.5\%$, $99\%$ thresholds). A formal incremental
test finds that augmenting a parsimonious volatility baseline with the full CTER conditioning set does
\emph{not} improve out-of-sample episode-severity forecasts---indeed it is significantly worse at the
$95\%$ and $97.5\%$ thresholds and statistically indistinguishable at $99\%$. The evidence therefore
supports a narrower conclusion: TER is a distinct path-dependent tail-risk \emph{functional} that captures
the temporal organization of extreme losses, while the \emph{forecasting} value of its conditional
extension (CTER) remains unresolved.

\medskip
\noindent\textbf{Keywords:} Tail risk; Value-at-Risk; Expected Shortfall; extreme-value theory; clustered
extremes; path-dependent risk; financial risk measurement; Conditional Tail Episode Risk; volatility
clustering.

\noindent\textbf{JEL Classification:} C22; C53; G17; G32.
\end{abstract}

\section{Introduction}\label{sec:intro}

Financial risk measurement traditionally focuses on the marginal distribution of portfolio or asset
losses. The two most widely used measures are Value-at-Risk and Expected Shortfall: VaR identifies a loss
quantile at a specified confidence level, while ES measures the expected loss beyond that quantile. Their
importance reflects both their statistical interpretation and their role in risk management, capital
allocation, and regulation.

Financial losses, however, are not independent draws from an unchanging marginal distribution. Extreme
losses cluster, volatility varies over time, and market stress generates \emph{sequences} of losses rather
than isolated observations. Consequently, two future loss paths can share an identical marginal
distribution---and therefore identical VaR and ES---while exhibiting very different temporal structure.
Consider
\[
A=(4,0,4,0,4,0), \qquad B=(4,4,4,0,0,0),
\]
with a tail threshold $q=2$. Both contain three losses of magnitude four and three zeros, so their marginal
distributions, and hence any marginal VaR or ES, coincide. Yet the losses in $A$ occur as three isolated
exceedances, whereas those in $B$ form one contiguous extreme episode. This distinction motivates the
central question:

\begin{quote}
\emph{Can financial tail risk be characterized by the temporal organization and cumulative severity of
extreme-loss episodes, rather than only by the marginal magnitude of individual tail observations?}
\end{quote}

To investigate it we introduce \textbf{Tail Episode Risk (TER)}. The idea is deliberately related to
established extreme-value theory: aggregate excess severity---the sum of exceedances over a high
threshold---has been studied for decades, and \citet{anderson1994} explicitly proposed aggregate excess as
a measure of extreme-event severity, emphasizing that severity depends not only on peak magnitude but also
on duration and temporal shape. The contribution here is therefore \emph{not} the invention of threshold
exceedances or aggregate excess. Rather, we formulate a financial risk functional around the upper tail of
the \emph{maximum cumulative threshold-exceedance episode} over a forecast horizon:
\[
\TER_{\alpha,\beta,H}
= Q_{\beta}\!\left( \max_{j\in\mathcal E_{H}} \sum_{k\in j}(L_k-q_\alpha)_+ \right),
\]
where $L_k$ is future loss, $q_\alpha$ a tail threshold, $\mathcal E_H$ the set of contiguous
threshold-exceedance episodes over horizon $H$, and $Q_\beta$ an upper-tail quantile of the resulting
episode-severity distribution. We then introduce \textbf{Conditional Tail Episode Risk (CTER)},
\[
\CTER_{\alpha,\beta,H,t}
= \Prob\!\left(O_{t,H}=1 \mid X_t\right)\,
  \ES_\beta\!\left[ S^{*}_{t,H} \mid O_{t,H}=1, X_t, Z_t \right],
\]
which separates two questions: how likely is the market to \emph{enter} an extreme-loss episode, and,
conditional on entering, how severe could the episode become?

We pose two distinct research questions.
\begin{description}
\item[RQ1.] Does TER capture temporal information about extreme losses that marginal VaR and ES do not?
\item[RQ2.] Does conditional forecasting of TER through CTER provide incremental predictive information
beyond conventional tail-risk models?
\end{description}
The distinction is central: a risk functional can encode information absent from VaR/ES without
necessarily producing superior forecasts in a finite real-world sample.

\paragraph{Contributions.} \emph{First}, we formulate TER as a path-dependent tail-risk functional based
on maximum cumulative threshold-exceedance severity. \emph{Second}, we introduce CTER as a conditional
probability--severity decomposition of episode risk. \emph{Third}, we provide theoretical and
simulation-based evidence that TER reacts to temporal clustering even when marginal tail behavior is held
approximately fixed. \emph{Fourth}, we provide an empirical evaluation on 2000--2026 data that reports both
positive and negative results and, in particular, does \emph{not} find that CTER universally improves
conventional VaR/ES-style forecasting.

\paragraph{Road map.} Section~\ref{sec:lit} reviews related risk measures. Section~\ref{sec:ter} defines
TER and Section~\ref{sec:cter} defines CTER. Section~\ref{sec:theory} establishes theoretical properties.
Section~\ref{sec:mc} presents the Monte Carlo study and Section~\ref{sec:stress} the adversarial stress
tests. Section~\ref{sec:design} describes the empirical design and the fully specified estimation pipeline.
Section~\ref{sec:bench} reports conventional VaR/ES benchmark results; Section~\ref{sec:cterresults}
reports CTER out-of-sample results; and Section~\ref{sec:incremental} reports the incremental-information
test. Section~\ref{sec:interpret} interprets the gap between measure innovation and forecasting
performance, Section~\ref{sec:relation} positions TER among existing measures, Section~\ref{sec:econ}
discusses economic and governance implications, Section~\ref{sec:limits} states limitations,
Section~\ref{sec:future} outlines future research, and Section~\ref{sec:conclusion} concludes.

\section{Related Literature}\label{sec:lit}

\paragraph{Value-at-Risk and Expected Shortfall.} For a loss $L$, $\VaR_\alpha(L)=Q_\alpha(L)$ and
$\ES_\alpha(L)=\E[L\mid L\ge \VaR_\alpha(L)]$ under a continuous distribution. VaR provides a threshold but
not the average magnitude beyond it; ES addresses this by averaging the tail. Modern forecast evaluation
requires care: \citet{fzg2016} and \citet{fz2016} established the joint elicitability of VaR and ES,
providing a foundation for consistent joint scoring and comparative backtesting
\citep{noldeziegel2017,pzc2019}. We evaluate conventional VaR/ES forecasts with a strictly consistent
Fissler--Ziegel scoring function (Section~\ref{sec:bench}); CTER forecasts a different statistical object
and is evaluated separately.

\paragraph{Aggregate excess and extreme-value theory.} Instead of the single maximum, an extreme event can
be characterized by the total amount by which observations exceed a threshold. \citet{anderson1994}
proposed the aggregate excess measure $S=\sum_{j=1}^{N}\zeta_j$, where $\zeta_j$ is an excess over a high
threshold, and emphasized that severity depends on duration and temporal shape. This literature establishes
a clear prior-art boundary: the sum of threshold exceedances is \emph{not} claimed here as an original
primitive. The proposed distinction is the financial functional $\max_j\sum_{k\in j}(L_k-q_\alpha)_+$
followed by a second tail operation over the resulting distribution.

\paragraph{Drawdown-based risk.} Maximum drawdown is another path-dependent concept; Conditional Expected
Drawdown (CED) formalizes the tail of maximum drawdown and is positively homogeneous, convex, and sensitive
to serial correlation \citep{goldbergmahmoud2017}. CED and TER are both path-dependent but have different
state variables:
\[
\mathrm{MDD}=\max_{u<v}(W_u-W_v)
\qquad\text{versus}\qquad
S^{*}_{t,H}=\max_{j}\sum_{k\in\mathcal K_{t,H,j}}(L_{t+k}-q_{t,\alpha})_+ ,
\]
where $W$ denotes wealth. TER depends on threshold exceedances; CED depends on peak-to-trough wealth
deterioration. TER is therefore not a drawdown measure in disguise.

\paragraph{Dynamic extreme-value models.} A large literature develops dynamic EVT models for financial risk
forecasting. \citet{candiaherrera2024} review dynamic threshold-exceedance models and find that
marked-point-process and time-series approaches can perform strongly, while no single specification
dominates. Recent work continues to develop conditional extreme-value models with time-varying scale and
tail parameters \citep{lihouzhang2026} and dynamic tail-risk networks and mixed-frequency frameworks
\citep{zhangyanshen2026,wang2026jof}. These models improve estimation and forecasting of conventional tail
measures; CTER instead targets future extreme-episode occurrence and severity.

\section{Tail Episode Risk}\label{sec:ter}

\subsection{Loss path and exceedances}
Let the future loss path over horizon $H$ be $\mathbf L_{t,H}=(L_{t+1},\dots,L_{t+H})$ and let $q_{t,\alpha}$
denote a tail threshold. Define the (non-negative) exceedance and the exceedance indicator
\[
y_{t+k}=(L_{t+k}-q_{t,\alpha})_+,\qquad I_{t+k}=\ind\{L_{t+k}>q_{t,\alpha}\}.
\]
A \emph{tail episode} is a maximal contiguous run of indices with $I_{t+k}=1$; let $\mathcal E_{t,H}$ denote
the set of such episodes.\footnote{We write the exceedance \emph{value} as $y_{t+k}$ and reserve
$O_{t,H}$ (Section~\ref{sec:cter}) for the episode-\emph{occurrence} event, avoiding the symbol reuse that
would otherwise conflate a magnitude with an event.}

\subsection{Episode severity and TER}
For episode $j\in\mathcal E_{t,H}$, let $\mathcal K_{t,H,j}\subseteq\{1,\dots,H\}$ be its (contiguous) index
set and define its cumulative excess severity
\[
S_{t,H,j}=\sum_{k\in\mathcal K_{t,H,j}}(L_{t+k}-q_{t,\alpha})_+ ,
\]
and the worst-episode severity over the horizon
\[
S^{*}_{t,H}=\max_{j\in\mathcal E_{t,H}} S_{t,H,j}
\qquad(\text{with }S^{*}_{t,H}=0\text{ if no exceedance occurs}).
\]
The maximum operation focuses the measure on the single most severe contiguous episode. Tail Episode Risk
is the upper-tail quantile of this worst-episode severity, conditional on information $\mathcal F_t$:
\[
\boxed{\;\TER_{\alpha,\beta,H,t}=Q_\beta\!\left(S^{*}_{t,H}\mid\mathcal F_t\right)\;}
\]
where $\alpha$ sets the threshold for an extreme observation, $\beta$ the tail level applied to the
distribution of maximum episode severity, and $\mathcal F_t$ the information set at date $t$. The two
probability levels have distinct roles and need not be equal.

\section{Conditional Tail Episode Risk}\label{sec:cter}

\subsection{Episode probability}
Define the episode-occurrence event
\[
O_{t,H}=\ind\!\left\{\max_{1\le k\le H}L_{t+k}>q_{t,\alpha}\right\},
\qquad p_t=\Prob(O_{t,H}=1\mid X_t),
\]
where $X_t$ collects observable state information (lagged returns, realized volatility, volatility
acceleration, drawdown, cross-asset returns, and financial-stress indicators; see
Section~\ref{sec:design}).

\subsection{Conditional episode severity}
Let $Z_t$ contain episode-state information (recent exceedance frequency, current run length, recent
episode severity). Define the conditional upper-tail severity
\[
m_t=\ES_\beta\!\left[S^{*}_{t,H}\mid O_{t,H}=1, X_t, Z_t\right].
\]

\subsection{Definition of CTER}
\[
\boxed{\;\CTER_{\alpha,\beta,H,t}=p_t\,m_t
=\Prob(O_{t,H}=1\mid X_t)\,\ES_\beta\!\left[S^{*}_{t,H}\mid O_{t,H}=1,X_t,Z_t\right].\;}
\]
CTER thus factors episode risk into an \emph{episode probability} and a \emph{conditional episode
severity}. This is economically different from forecasting another VaR or ES: it explicitly separates the
likelihood of entering a stressed state from its potential severity.

\section{Theoretical Properties}\label{sec:theory}

\begin{proposition}[Monotonicity]
If $L'_k\ge L_k$ for all $k$ under a common threshold $q$, then $(L'_k-q)_+\ge(L_k-q)_+$, so no episode's
cumulative excess can decrease and $S^{*}(\mathbf L')\ge S^{*}(\mathbf L)$. Hence
$\TER(\mathbf L')\ge \TER(\mathbf L)$.
\end{proposition}

\begin{proposition}[Positive homogeneity under consistent scaling]
For $c>0$, if $L'_k=cL_k$ and $q'=cq$ then $(L'_k-q')_+=c(L_k-q)_+$, so $S^{*\prime}=cS^{*}$ and
$\TER(cL)=c\,\TER(L)$.
\end{proposition}

\begin{proposition}[Translation invariance under a translation-equivariant threshold rule]
If $L'_k=L_k+c$ and the threshold rule is translation-equivariant, $q'=q+c$, then $L'_k-q'=L_k-q$, so
$S^{*\prime}=S^{*}$ and $\TER'=\TER$.
\end{proposition}
\noindent A clarification is warranted, because VaR and ES are translation \emph{equivariant}
($\rho(L+c)=\rho(L)+c$) as monetary risk measures. TER is not a monetary risk measure but a severity
aggregate, and the natural analogue is a \emph{deviation} functional, which satisfies translation
\emph{invariance} ($D(L+c)=D(L)$). Proposition~3 is exactly this: when the threshold shifts with the losses,
the exceedances---and hence $S^{*}$---are unchanged, so TER is invariant, not equivariant. TER neither
satisfies nor targets the monetary equivariance property; it measures the magnitude of a
threshold-exceedance episode rather than a capital requirement. The qualifications also matter for
homogeneity: both properties hold \emph{conditional} on the threshold being scaled or shifted jointly with
the losses, and because $q_{t,\alpha}$ is in practice estimated dynamically, they are properties of the
joint (loss, threshold) rule rather than unconditional axioms.

\begin{proposition}[Temporal-clustering sensitivity]
For $A=(4,0,4,0,4,0)$ and $B=(4,4,4,0,0,0)$ with $q=2$, the three exceedances in $A$ are separate
singletons, giving $S^{*}(A)=2$, whereas in $B$ they form one contiguous episode, giving $S^{*}(B)=6$.
Since the two paths are marginal permutations of one another, $\VaR(A)=\VaR(B)$ and $\ES(A)=\ES(B)$, yet
$S^{*}(A)\ne S^{*}(B)$ and therefore $\TER(A)\ne\TER(B)$.
\end{proposition}
\noindent This is the central conceptual distinction; Figure~\ref{fig:ab} illustrates it. (On a single path
$S^{*}$ is a scalar; $\TER=Q_\beta(S^{*})$ becomes a genuine quantile once $S^{*}$ is viewed as a random
variable over the horizon distribution.)

\begin{figure}[t]
\centering
\includegraphics[width=0.72\textwidth]{figs/ab.pdf}
\caption{Two loss paths with identical marginals (three losses of $4$, three of $0$) and hence identical VaR
and ES, but different worst-episode severity $S^{*}$ at threshold $q=2$.}
\label{fig:ab}
\end{figure}

\paragraph{Coherence.} We do \emph{not} claim that TER satisfies the full coherence axioms. Subadditivity
in particular requires separate analysis because the threshold is distribution-dependent and because
episode formation can change under portfolio aggregation. TER should therefore be treated as a
path-dependent tail-risk \emph{functional}, not automatically as a coherent risk measure.

\section{Monte Carlo Simulation}\label{sec:mc}

\paragraph{Objective.} We test whether TER is sensitive to temporal dependence when the marginal
innovation distribution is held approximately fixed. We simulate stationary AR(1) processes
$x_t=\rho x_{t-1}+\sqrt{1-\rho^2}\,\varepsilon_t$, so the marginal variance is invariant to $\rho$, and
compute VaR/ES on the marginal together with $Q_\beta(S^{*})$ over rolling $H=10$ windows
($20{,}000$ windows per configuration; threshold $\alpha=95\%$). VaR and ES are computed from the stationary
marginal distribution implied by the standardized AR(1) process, whereas TER is computed from the rolling
horizon paths; this isolates the dependence channel, since the marginal is held fixed by construction while
only $\rho$ varies.

\begin{table}[t]
\centering
\caption{Gaussian AR(1) dependence experiment ($\varepsilon\sim N(0,1)$, standardized).}
\label{tab:mc-normal}
\begin{tabular}{rrrrr}
\toprule
$\rho$ & VaR $95\%$ & ES $95\%$ & TER $95\%$ & TER $99\%$ \\
\midrule
0.0 & 1.633 & 2.059 & 0.981 & 1.508 \\
0.3 & 1.645 & 2.053 & 1.010 & 1.737 \\
0.7 & 1.655 & 2.046 & 1.065 & 2.580 \\
0.9 & 1.668 & 2.093 & 1.060 & 4.541 \\
\bottomrule
\end{tabular}
\end{table}

\begin{table}[t]
\centering
\caption{Student-$t$ AR(1) dependence experiment ($\varepsilon\sim t_5$, standardized to unit variance).}
\label{tab:mc-t}
\begin{tabular}{rrrrr}
\toprule
$\rho$ & VaR $95\%$ & ES $95\%$ & TER $95\%$ & TER $99\%$ \\
\midrule
0.0 & 1.560 & 2.233 & 1.494 & 3.130 \\
0.3 & 1.562 & 2.199 & 1.530 & 2.979 \\
0.7 & 1.579 & 2.134 & 1.482 & 4.141 \\
0.9 & 1.622 & 2.106 & 1.304 & 5.281 \\
\bottomrule
\end{tabular}
\end{table}

Tables~\ref{tab:mc-normal}--\ref{tab:mc-t} and Figure~\ref{fig:mc} report the result. In the Gaussian case,
$95\%$ VaR and ES are essentially flat in $\rho$ (near the standard-normal values $1.645$ and $2.063$),
while $\TER_{99}$ rises from $\approx 1.5$ to $\approx 4.5$. The Student-$t$ case, with innovations
standardized to unit variance so that the marginal is held fixed, shows the same qualitative pattern
($\TER_{99}$ rising from $3.1$ to $5.3$).\footnote{Standardizing the $t$ innovations is important: an
\emph{un}-standardized $t_5$ has variance $5/3$, which by itself inflates all tail measures and confounds
the dependence channel. With standardization, the residual rise in $\TER_{99}$ is attributable to
clustering.} We summarize:
\[
\boxed{\text{Temporal dependence can materially raise episode risk without proportionally raising marginal
VaR/ES.}}
\]
This demonstrates that TER measures an \emph{additional} dimension of path information; it does not by
itself demonstrate a superior forecast.

\begin{figure}[t]
\centering
\includegraphics[width=0.62\textwidth]{figs/mc.pdf}
\caption{Gaussian AR(1) experiment. As serial dependence $\rho$ increases, VaR and ES are flat while TER
(especially at the $99\%$ level) rises sharply.}
\label{fig:mc}
\end{figure}

\section{Adversarial Stress Tests}\label{sec:stress}

To probe behavior directly, we construct six $60$-day loss paths under a \emph{common, fixed} threshold
$q=2.5$ (in standardized loss units), so that differences in the measures reflect temporal organization
rather than a path-specific threshold. Table~\ref{tab:stress} reports $95\%$ VaR, $95\%$ ES, and the
worst-episode severity $S^{*}$.

\begin{table}[t]
\centering
\caption{Adversarial stress tests (fixed threshold $q=2.5$). $S^{*}$ is the worst contiguous
episode severity.}
\label{tab:stress}
\begin{tabular}{lrrr}
\toprule
Scenario & VaR $95\%$ & ES $95\%$ & $S^{*}$ \\
\midrule
Isolated crash                       & 1.42 & 4.00 & 6.50 \\
Tail thickening (separated spikes)   & 6.00 & 6.00 & 3.50 \\
Liquidity collapse (sustained run)   & 3.50 & 3.50 & 12.00 \\
Volatility explosion (cluster)       & 6.42 & 7.68 & 25.62 \\
Broad contagion (all elevated)       & 3.39 & 4.01 & 2.40 \\
Escalating deterioration             & 3.58 & 4.52 & 7.80 \\
\bottomrule
\end{tabular}
\end{table}

The decisive comparison is tail thickening versus liquidity collapse. Tail thickening has the
\emph{highest} VaR and ES ($6.00$, $6.00$) because it contains several large \emph{isolated} observations,
yet its worst episode is small ($S^{*}=3.5$). The liquidity collapse has \emph{lower} VaR and ES ($3.50$)
but the sustained run yields a much larger $S^{*}=12.0$. Under these constructed paths, VaR and ES assign
larger marginal-tail severity to tail thickening, whereas TER assigns larger episode severity to the
sustained liquidity-collapse path; the measures therefore distinguish different dimensions of stress rather
than producing a common risk ordering. The cross-asset (``broad contagion'') scenario also exposes a
limitation: the present TER is univariate and a
complete contagion measure would require an explicit multivariate extension.

\section{Empirical Design}\label{sec:design}

\subsection{Data}
We use daily log returns for six assets spanning 2000--2026: the S\&P~500 and Nasdaq indices, the iShares
20+ Year Treasury Bond ETF (TLT), EUR/USD, USD/JPY, and Bitcoin. Data availability differs by asset: the
equity indices and the two currency pairs begin in 2000, TLT in July 2002, and Bitcoin in September 2014.
State variables are the VIX index (daily), the Chicago Fed National Financial Conditions Index (NFCI), and
the St.\ Louis Fed Financial Stress Index (STLFSI); the two Fed indices are weekly and are carried forward
to daily frequency by last-observation-carried-forward (LOCF), so that only information available up to
date~$t$ enters $X_t$. The three stress series are used from their per-series source files rather than from
a merged file, so that the weekly indices are never treated as daily observations.

\paragraph{Calendar convention.} Because the paper studies temporal clustering, the observation calendar is
stated explicitly. For the exchange-traded assets, observations follow the relevant market calendar; Bitcoin
trades continuously and is observed on a calendar-day basis. The horizon is therefore defined in
\emph{observations} rather than calendar time: a 10-day observation horizon is ten trading days for the
exchange-traded assets and ten calendar days for Bitcoin. All quantities ($q_{t,\alpha}$, episodes,
$S^{*}_{t,H}$) are defined on each asset's own observation index.

\paragraph{Windows and horizon.} The rolling estimation window is $W=1{,}250$ observations, the forecast
horizon is $H=10$ observations, and the threshold levels are $\alpha\in\{95\%,97.5\%,99\%\}$. Because the
threshold window and a $252$-observation training burn-in consume the first $1{,}502$ observations,
out-of-sample forecasting begins at an \emph{asset-specific} date: December~2005 for the S\&P~500, Nasdaq,
EUR/USD, and USD/JPY; July~2008 for TLT; and October~2018 for Bitcoin (reflecting its later inception and
continuous calendar). The out-of-sample counts are $4{,}827$ (equity indices), $4{,}802$ (currencies),
$4{,}183$ (TLT), and $2{,}496$ (Bitcoin).

\subsection{Threshold and realized targets}
At each date $t$, $q_{t,\alpha}$ is the trailing-window empirical $\alpha$-quantile of losses $L=-r$. Over
the next $H$ days we compute the exceedances, form contiguous episodes, and record the realized worst-episode
severity $S^{*}_{t,H}$ and occurrence $O_{t,H}=\ind\{S^{*}_{t,H}>0\}$.

\subsection{Conditioning information ($X_t,Z_t$)}
All features use information available through date $t$. The state vector $X_t$ contains: today's loss;
trailing $5$- and $10$-day cumulative losses; $20$- and $60$-day realized volatility; volatility
acceleration ($\mathrm{rv}_{20}/\mathrm{rv}_{60}$); $60$-day drawdown; the VIX level and its $5$-day change;
NFCI; STLFSI; and the S\&P~500 trailing $5$-day loss as a systemic proxy. The episode-state vector $Z_t$
adds: the trailing $60$-day exceedance frequency relative to $q_{t,\alpha}$, the current in-episode run
length, and the recent within-window episode severity. The full specification is given in
Appendix~\ref{app:repro}.

\subsection{Estimators and evaluation}
The framework is deliberately estimator-agnostic; we fix one transparent specification. Episode occurrence
$p_t$ is estimated with a gradient-boosted classifier and episode severity with a gradient-boosted
regressor (histogram GBMs; $200$ trees, depth $3$, learning rate $0.06$, $\ell_2$ regularization $1.0$).
Models are fit by expanding-window walk-forward with annual\,($378$-day) refits and a purge gap of $H$ days
between the last training label and the first test date, eliminating look-ahead from the horizon-$H$
targets. Occurrence is evaluated by Brier score and AUC; severity by MAE and RMSE against realized
$S^{*}_{t,H}$; conventional VaR/ES benchmarks by a strictly consistent Fissler--Ziegel score
(Section~\ref{sec:bench}). This avoids scoring different statistical targets with a single function.

\section{Conventional VaR/ES Benchmark Results}\label{sec:bench}

We evaluate four one-day-ahead VaR/ES forecasters---Historical Simulation, an EWMA (RiskMetrics,
$\lambda=0.94$) Gaussian model, a peaks-over-threshold EVT/GPD model, and a GARCH(1,1)-$t$ model---using the
$0$-homogeneous Fissler--Ziegel loss \citep{fz2016,pzc2019}. For returns $r$, forecast VaR $v$ and ES $e$
(both negative in the lower tail) at tail probability $\alpha$,
\begin{equation}
\mathrm{FZ0}(r,v,e)= -\frac{1}{\alpha e}\,\ind\{r\le v\}\,(v-r)\;+\;\frac{v}{e}\;+\;\ln(-e)\;-\;1,
\label{eq:fz0}
\end{equation}
which is strictly consistent for the (VaR, ES) pair; lower average scores are better. Table~\ref{tab:fz}
reports the mean score across the six assets and the associated breach rates.

\begin{table}[t]
\centering
\caption{Mean FZ0 score across assets (lower is better) and mean breach rate. Best model per level in bold.}
\label{tab:fz}
\begin{tabular}{lrrrr}
\toprule
& Hist.\ Sim. & EWMA-Normal & EVT-POT & GARCH-$t$ \\
\midrule
\multicolumn{5}{l}{\emph{Mean FZ0 score}}\\
\quad $95\%$   & $-3.627$ & $-3.768$ & $-3.618$ & $\mathbf{-3.789}$ \\
\quad $97.5\%$ & $-3.395$ & $-3.535$ & $-3.383$ & $\mathbf{-3.590}$ \\
\quad $99\%$   & $-3.127$ & $-3.177$ & $-3.092$ & $\mathbf{-3.350}$ \\
\midrule
\multicolumn{5}{l}{\emph{Mean breach rate} (target $=1-\alpha$)}\\
\quad $95\%$   & 0.0500 & 0.0548 & 0.0497 & 0.0531 \\
\quad $97.5\%$ & 0.0259 & 0.0339 & 0.0258 & 0.0290 \\
\quad $99\%$   & 0.0122 & 0.0199 & 0.0115 & 0.0124 \\
\bottomrule
\end{tabular}
\end{table}

GARCH-$t$ attains the best FZ0 score at every level, with EWMA-Normal a strong second; Historical Simulation
and EVT are best calibrated by breach rate but weaker under the joint score. Notably, the EWMA-Normal model
over-breaches at the $99\%$ level ($1.99\%$ versus a $1\%$ target). We therefore use GARCH-$t$ (strongest
under FZ0) and EWMA (strong and simple) as the competitive baselines against which any incremental CTER
value must be judged---rather than a deliberately weak model that would flatter the new framework.

\section{CTER Out-of-Sample Results}\label{sec:cterresults}

Table~\ref{tab:cter} reports occurrence discrimination (Brier, AUC) and severity accuracy (MAE, RMSE) by
asset and threshold. Two patterns stand out. Severity forecasts exhibit measurable cross-market
predictability---their scale tracks each market's volatility (MAE is smallest for FX, largest for
Bitcoin)---although, as Section~\ref{sec:incremental} shows, the full CTER conditioning set does not improve
on a parsimonious baseline. Occurrence discrimination, by contrast, is weak and \emph{decays into the tail}:
the cross-asset mean AUC falls from $0.56$ ($95\%$) to $0.54$ ($97.5\%$) to $0.47$ ($99\%$)---below random at
the deepest level (Figure~\ref{fig:auc}). The bond market (TLT) is the exception, with AUC $\approx0.64$--$0.67$
at all levels, indicating that fixed-income episodes are more predictable than equity or crypto episodes. The
two components of the decomposition therefore behave asymmetrically: the probability component is weak, while
the severity target carries more predictable structure. The decomposition is thus conceptually useful even
though the current conditioning set does not reliably forecast episode occurrence, and CTER should not be
presented as a universal early-warning system.

\begin{table}[t]
\centering
\caption{CTER out-of-sample results by asset and threshold. ``OOS start'' is the first out-of-sample date
for that asset (windows differ by data availability). Base rate is the realized episode frequency; AUC below
$0.5$ indicates worse-than-random discrimination.}
\label{tab:cter}
\begin{tabular}{llrrrrr}
\toprule
Asset (OOS start) & Level & Base rate & Brier & AUC & MAE & RMSE \\
\midrule
S\&P 500 (Dec 2005) & $95\%$   & 0.317 & 0.249 & 0.563 & 0.00652 & 0.01436 \\
                    & $97.5\%$ & 0.179 & 0.209 & 0.471 & 0.00489 & 0.01256 \\
                    & $99\%$   & 0.098 & 0.130 & 0.405 & 0.00318 & 0.01012 \\
Nasdaq (Dec 2005)   & $95\%$   & 0.328 & 0.257 & 0.563 & 0.00661 & 0.01415 \\
                    & $97.5\%$ & 0.193 & 0.182 & 0.519 & 0.00470 & 0.01280 \\
                    & $99\%$   & 0.102 & 0.114 & 0.423 & 0.00306 & 0.00977 \\
TLT (Jul 2008)      & $95\%$   & 0.371 & 0.227 & 0.672 & 0.00304 & 0.00637 \\
                    & $97.5\%$ & 0.200 & 0.162 & 0.664 & 0.00206 & 0.00570 \\
                    & $99\%$   & 0.093 & 0.094 & 0.644 & 0.00134 & 0.00520 \\
EUR/USD (Dec 2005)  & $95\%$   & 0.337 & 0.241 & 0.597 & 0.00196 & 0.00336 \\
                    & $97.5\%$ & 0.186 & 0.176 & 0.601 & 0.00117 & 0.00239 \\
                    & $99\%$   & 0.090 & 0.096 & 0.467 & 0.00055 & 0.00153 \\
USD/JPY (Dec 2005)  & $95\%$   & 0.372 & 0.258 & 0.599 & 0.00392 & 0.00720 \\
                    & $97.5\%$ & 0.202 & 0.186 & 0.581 & 0.00296 & 0.00591 \\
                    & $99\%$   & 0.119 & 0.132 & 0.473 & 0.00205 & 0.00504 \\
Bitcoin (Oct 2018)  & $95\%$   & 0.251 & 0.308 & 0.377 & 0.02858 & 0.04827 \\
                    & $97.5\%$ & 0.130 & 0.191 & 0.427 & 0.01920 & 0.03858 \\
                    & $99\%$   & 0.054 & 0.131 & 0.403 & 0.00937 & 0.02754 \\
\bottomrule
\end{tabular}
\end{table}

\begin{figure}[t]
\centering
\includegraphics[width=0.72\textwidth]{figs/auc.pdf}
\caption{Episode-occurrence AUC by asset and threshold. Discrimination is weak overall and, except for
bonds (TLT), falls toward or below $0.5$ at the $99\%$ threshold.}
\label{fig:auc}
\end{figure}

As an economic reality check, Figure~\ref{fig:episodes} plots the realized $10$-day episode severity for the
S\&P~500 over 2000--2026. The measure isolates precisely the sustained crises one would want it to flag: its
worst distinct episodes are March~2020 (COVID, $S^{*}=11.3\%$), October--November~2008 (GFC), the
March~2025 selloff, August~2011, and the 2002 and 2000 dot-com unwinds. This confirms that $S^{*}$ is
economically meaningful even where its \emph{conditional} forecast is difficult.

\begin{figure}[t]
\centering
\includegraphics[width=0.92\textwidth]{figs/episodes.pdf}
\caption{Realized S\&P~500 Tail Episode severity $S^{*}_{t,H}$ ($H=10$, threshold $q$ at the full-sample
$95\%$ loss quantile, $1.85\%$). Labeled peaks are the worst distinct episodes.}
\label{fig:episodes}
\end{figure}

\section{Incremental Information Test}\label{sec:incremental}

The decisive question (RQ2) is whether CTER adds information beyond a competitive conventional model. We
compare two forecasts of the \emph{same} target $S^{*}_{t,H}$, produced by the \emph{same} gradient-boosted
regressor under the \emph{same} walk-forward protocol, differing only in their inputs. The \emph{baseline}
regressor receives three EWMA-derived inputs---the EWMA volatility $\hat\sigma^{\mathrm{EWMA}}_t$ and its
induced one-day VaR and ES, $\hat\sigma^{\mathrm{EWMA}}_t\Phi^{-1}(\alpha)$ and
$-\hat\sigma^{\mathrm{EWMA}}_t\phi(\Phi^{-1}(\alpha))/\alpha$. The \emph{augmented} regressor receives these
three inputs \emph{plus} the full $X_t,Z_t$ conditioning set. The map from EWMA inputs to the 10-observation
episode-severity forecast is thus the gradient-boosted regressor itself; no closed-form scaling is imposed,
which is why both models must share the same estimator so that the comparison isolates the incremental value
of the conditioning information. Let
$\Delta_t=|S^{*}_t-\hat S^{*}_{\text{aug},t}|-|S^{*}_t-\hat S^{*}_{\text{base},t}|$; negative values favor
CTER. We aggregate equal-weight across assets and report the mean $\Delta$, a moving-block bootstrap $95\%$
confidence interval (block length $20$), and a Diebold--Mariano $p$-value with a Newey--West HAC variance
(lag $H-1=9$).

\begin{table}[t]
\centering
\caption{Incremental-information test. $\Delta=$ MAE(augmented) $-$ MAE(baseline); \emph{negative} favors
CTER. Positive, significant values indicate the conditioning set does not help.}
\label{tab:incremental}
\begin{tabular}{lrrr}
\toprule
Level & Mean $\Delta$ & $95\%$ block-bootstrap CI & DM $p$-value \\
\midrule
$95\%$   & $+0.001231$ & $[+0.000729,\ +0.001790]$ & $<0.001$ \\
$97.5\%$ & $+0.000722$ & $[+0.000339,\ +0.001141]$ & $<0.001$ \\
$99\%$   & $+0.000170$ & $[-0.000076,\ +0.000409]$ & $0.077$ \\
\bottomrule
\end{tabular}
\end{table}

The result (Table~\ref{tab:incremental} and Figure~\ref{fig:incremental}) is a clear negative. At the $95\%$
and $97.5\%$ thresholds the augmented model is \emph{significantly worse} than the parsimonious volatility
baseline ($\Delta>0$, $p<0.001$), and at $99\%$ the difference is statistically insignificant (the CI
includes zero, $p=0.077$). This pattern is consistent with the additional conditioning variables introducing
estimation variance without sufficient incremental predictive signal in this specification and sample. The
appropriate conclusion is:
\begin{quote}
The 2000--2026 cross-asset sample provides \emph{no} evidence that CTER conditioning improves aggregate
out-of-sample episode-severity forecasting over a competitive volatility baseline, and some evidence that
it degrades it at the $95\%$ and $97.5\%$ thresholds.
\end{quote}

\begin{figure}[t]
\centering
\includegraphics[width=0.66\textwidth]{figs/incremental.pdf}
\caption{Incremental CTER forecast performance. Points are the mean MAE difference (augmented minus
baseline) with $95\%$ moving-block bootstrap intervals; values to the left of the dashed zero line would
favor CTER. The augmented model is significantly worse at $95\%$ and $97.5\%$ and indistinguishable from the
baseline at $99\%$.}
\label{fig:incremental}
\end{figure}

\section{Why the Distinction Matters}\label{sec:interpret}

The theoretical and empirical findings are not contradictory. The simulations (Section~\ref{sec:mc}) show
that $\TER$ captures the temporal organization of extreme losses; the forecasting experiments
(Sections~\ref{sec:cterresults}--\ref{sec:incremental}) show that CTER does not yet demonstrate incremental
predictive value. A risk measure can encode a meaningful dimension of risk while being hard to forecast.
The pipeline
\[
\text{Risk functional}\ \to\ \text{Target}\ \to\ \text{Estimator}\ \to\ \text{Forecast}\ \to\ \text{Decision}
\]
is supported here at the first two links and remains open at the later ones. Separating
\emph{risk-functional innovation} from \emph{forecasting performance} is the paper's central methodological
point.

\section{Relation to Existing Risk Measures}\label{sec:relation}

\begin{table}[h]
\centering
\caption{Primary question addressed by each measure.}
\label{tab:relation}
\begin{tabular}{ll}
\toprule
Measure & Primary question \\
\midrule
VaR & How large is the loss threshold? \\
ES & How severe is the marginal tail? \\
Aggregate excess & How much total exceedance occurs within an extreme event? \\
CED & How large can peak-to-trough drawdown become? \\
TER & How severe can the worst cumulative threshold-exceedance episode become? \\
CTER & Given current conditions, how likely and how severe is that episode? \\
\bottomrule
\end{tabular}
\end{table}
\noindent TER occupies a position between marginal tail risk and fully path-dependent stress analysis.
Table~\ref{tab:contribution} restates this positioning against the specific structural properties that
distinguish the proposed functionals from existing measures.

\begin{table}[h]
\centering
\caption{Contribution relative to existing measures. ``Path dependent'' indicates dependence on the ordering
of observations; ``temporal clustering'' indicates sensitivity to whether exceedances are contiguous.}
\label{tab:contribution}
\begin{tabular}{lllll}
\toprule
Measure & Object & Path dependent? & Temporal clustering? & Status \\
\midrule
VaR & Quantile of loss & No & No & Existing \\
ES & Mean tail loss & No & No & Existing \\
Aggregate excess & Total threshold exceedance & Event-level & Yes & Related (prior art) \\
CED & Maximum drawdown & Yes & Yes & Related \\
TER & Max cumulative threshold-exceedance episode & Yes & Yes & \textbf{Proposed} \\
CTER & Conditional prob.\ $\times$ tail episode severity & Yes & Yes & \textbf{Proposed} \\
\bottomrule
\end{tabular}
\end{table}

\section{Economic and Governance Implications}\label{sec:econ}

TER is not proposed as a replacement for regulatory VaR or ES but as a complementary monitoring layer. A
conventional system reporting, say, $\VaR=5\%$ and $\ES=7\%$ could additionally report an episode
probability and a conditional episode severity, communicating the likelihood and potential severity of a
\emph{sustained} extreme-loss episode. Because the framework makes the temporal structure of tail events
explicit, a risk system could separately track episode probability, severity, duration, escalation, and
clustering---more interpretable than a single opaque score. That additional structure, however, creates
additional model risk: the framework depends on threshold selection, episode definition, horizon, tail
level, probability calibration, severity estimation, and dependence assumptions, and should be subject to
model validation rather than treated as superior because it is more complex. The empirical results here
reinforce that caution: the conditional layer did not add forecasting value in this sample.

\section{Limitations}\label{sec:limits}

\emph{(i) Prior art.} The cumulative threshold-exceedance component is established in EVT
\citep{anderson1994}; the potential contribution is the financial maximum-episode functional and its
conditional probability--severity extension. \emph{(ii) No empirical superiority.} The incremental tests
are, at best, insignificant and, at worst, unfavorable to CTER. \emph{(iii) Univariate formulation.}
Cross-asset contagion requires an explicit multivariate extension. \emph{(iv) Threshold and horizon
sensitivity.} Different $q_\alpha$ and $H$ define different objects; stability warrants systematic study.
\emph{(v) Sparse targets.} Episode occurrence and severity are hard to forecast because extreme events are
rare, especially at $99\%$. \emph{(vi) Estimator scope.} A fully covariate-conditioned EVT specification for
$\ES_\beta(S^{*}\mid O,X_t,Z_t)$ remains unexplored and could differ from the gradient-boosted
implementation used here.

\section{Future Research}\label{sec:future}

Natural extensions include a multivariate CTER for synchronized cross-asset episodes; a liquidity-adjusted
variant incorporating spreads, depth, turnover, and price impact; an options-implied variant using implied
volatility, skew, and the variance risk premium; state-dependent EVT with $\xi_t=f(X_t)$, $\sigma_t=g(X_t)$
in a conditional GPD; richer machine-learning estimators (quantile forests, distributional networks,
temporal transformers) with the functional held fixed while only the estimator varies; and a
decision-theoretic evaluation asking whether CTER changes capital, position sizing, or risk limits rather
than relying only on statistical forecast error.

\section{Conclusion}\label{sec:conclusion}

Financial tail risk has at least two dimensions: the \emph{magnitude} of individual tail losses and the
\emph{temporal organization} of those losses. VaR and ES are built around the first; TER explicitly
incorporates the second. We showed constructively and by simulation that two paths can share identical
marginal VaR and ES while differing in TER, and that increasing serial dependence can sharply raise episode
risk while leaving marginal measures nearly unchanged. The empirical evidence is more cautious: on
2000--2026 data across six markets, conditional forecasting of episode risk through CTER did not improve
on---and at moderate thresholds underperformed---a competitive volatility baseline, and episode-occurrence
discrimination was weak in the deep tail. The evidence therefore does not justify presenting CTER as a
replacement for VaR or ES. It supports a narrower, more precise conclusion:
\begin{quote}
Tail Episode Risk provides a distinct path-dependent representation of extreme-loss severity that preserves
temporal information discarded by marginal tail measures. Conditional Tail Episode Risk provides a framework
for incorporating current-state information into that representation, but its empirical forecasting value in
real markets remains an open question.
\end{quote}
The broader implication is that risk measurement may benefit from moving beyond how large an individual tail
loss can become toward the complementary question of how extreme losses organize themselves through time.

\bigskip
\noindent\textbf{Reproducibility.} All figures and tables are generated from daily 2000--2026 data by a
single pipeline. The manuscript specifies the estimator, feature set, evaluation protocol, and configuration
required to reproduce the reported analysis from the underlying data (Section~\ref{sec:design} and
Appendix~C); the accompanying code repository contains the data-preprocessing, estimation, forecast, and
scoring code together with fixed random seeds.

\begin{thebibliography}{99}
\small
\bibitem[Anderson(1994)]{anderson1994} Anderson, C.\ W.\ (1994). The aggregate excess measure of severity
of extreme events. \emph{Journal of Research of the National Institute of Standards and Technology},
99(4), 555--561. \href{https://doi.org/10.6028/jres.099.054}{doi:10.6028/jres.099.054}.

\bibitem[Candia and Herrera(2024)]{candiaherrera2024} Candia, C., \& Herrera, R.\ (2024). An empirical
review of dynamic extreme value models for forecasting Value-at-Risk, Expected Shortfall and expectile.
\emph{Journal of Empirical Finance}, 77, 101488.
\href{https://doi.org/10.1016/j.jempfin.2024.101488}{doi:10.1016/j.jempfin.2024.101488}.

\bibitem[Fissler et al.(2016)]{fzg2016} Fissler, T., Ziegel, J.\ F., \& Gneiting, T.\ (2016). Expected
Shortfall is jointly elicitable with Value-at-Risk---implications for backtesting. \emph{Risk Magazine},
58--61.

\bibitem[Fissler and Ziegel(2016)]{fz2016} Fissler, T., \& Ziegel, J.\ F.\ (2016). Higher order
elicitability and Osband's principle. \emph{Annals of Statistics}, 44(4), 1680--1707.

\bibitem[Goldberg and Mahmoud(2017)]{goldbergmahmoud2017} Goldberg, L.\ R., \& Mahmoud, O.\ (2017).
Drawdown: from practice to theory and back again. \emph{Mathematics and Financial Economics}, 11, 275--297.
\href{https://doi.org/10.1007/s11579-016-0181-9}{doi:10.1007/s11579-016-0181-9}.

\bibitem[Li et al.(2026)]{lihouzhang2026} Li, J., Hou, W., \& Zhang, Z.\ (2026). A coupled autoregressive
extreme-value model for dynamic tail risk with risk spirals. \emph{Finance Research Letters}, 105, 110187.
\href{https://doi.org/10.1016/j.frl.2026.110187}{doi:10.1016/j.frl.2026.110187}.

\bibitem[Nolde and Ziegel(2017)]{noldeziegel2017} Nolde, N., \& Ziegel, J.\ F.\ (2017). Elicitability and
backtesting: perspectives for banking regulation. \emph{Annals of Applied Statistics}, 11(4), 1833--1874.

\bibitem[Patton et al.(2019)]{pzc2019} Patton, A.\ J., Ziegel, J.\ F., \& Chen, R.\ (2019). Dynamic
semiparametric models for expected shortfall (and Value-at-Risk). \emph{Journal of Econometrics}, 211(2),
388--413.

\bibitem[Zhang et al.(2026a)]{zhangyanshen2026} Zhang, D., Yan, X., \& Shen, F.\ (2026). Dynamic financial
tail risk networks: a backtesting-based conditional expected shortfall approach. \emph{PLOS ONE}, 21(6),
e0351966. \href{https://doi.org/10.1371/journal.pone.0351966}{doi:10.1371/journal.pone.0351966}.

\bibitem[Wang et al.(2026)]{wang2026jof} Wang, S., Wang, J., Lv, M., \& Zhang, D.\ (2026). Forecasting
financial risk with minute-level transaction data and economic policy uncertainty: a new mixed-frequency
time-varying forecasting framework. \emph{Journal of Forecasting}, 45(4), 1681--1713.
\href{https://doi.org/10.1002/for.70101}{doi:10.1002/for.70101}.
\end{thebibliography}

\appendix
\section{Core Definitions}\label{app:defs}
\begin{align*}
\text{Tail episode set:}\quad & \mathcal E_{t,H}=\{\,j:\ L_{t+k}>q_{t,\alpha}\ \text{for contiguous }k\,\}.\\
\text{Episode severity:}\quad & S_{t,H,j}=\sum_{k\in\mathcal K_{t,H,j}}(L_{t+k}-q_{t,\alpha})_+.\\
\text{Maximum episode severity:}\quad & S^{*}_{t,H}=\max_{j\in\mathcal E_{t,H}}S_{t,H,j}.\\
\text{Tail Episode Risk:}\quad & \TER_{\alpha,\beta,H,t}=Q_\beta(S^{*}_{t,H}\mid\mathcal F_t).\\
\text{Conditional TER:}\quad & \CTER_{\alpha,\beta,H,t}=\Prob(O_{t,H}=1\mid X_t)\,
\ES_\beta[S^{*}_{t,H}\mid O_{t,H}=1,X_t,Z_t].
\end{align*}

\section{Empirical Evaluation Protocol}\label{app:protocol}
Daily data, 2000--2026, with asset-specific out-of-sample start dates (2005 for equities and FX, 2008 for
TLT, 2018 for Bitcoin); $W=1{,}250$-observation rolling threshold window; $H=10$-observation
horizon; levels $95\%,97.5\%,99\%$; assets S\&P~500, Nasdaq, TLT, EUR/USD, USD/JPY, Bitcoin. Conventional
VaR/ES models are scored with the FZ0 loss \eqref{eq:fz0}. CTER occurrence is scored by Brier and AUC and
severity by MAE and RMSE. Incremental value compares an EWMA-only baseline with a fully conditioned model on
the same target with paired, block-bootstrapped, HAC-adjusted tests. No post-out-of-sample tuning is used.

\section{Reproducibility Details}\label{app:repro}
\textbf{Features $X_t$:} today's loss $L_t$; trailing $5$- and $10$-day cumulative loss; $20$- and $60$-day
realized volatility; volatility acceleration $\mathrm{rv}_{20}/\mathrm{rv}_{60}$; $60$-day drawdown; VIX
level; VIX $5$-day change; NFCI; STLFSI; S\&P~500 trailing $5$-day loss. \textbf{Episode state $Z_t$:}
trailing $60$-day exceedance frequency relative to $q_{t,\alpha}$; current in-episode run length; recent
within-window episode severity. \textbf{Estimators:} histogram gradient-boosting classifier (occurrence)
and regressor (severity), $200$ trees, max depth $3$, learning rate $0.06$, $\ell_2$ regularization $1.0$.
\textbf{Walk-forward:} expanding window, first fit after $W+252$ observations, refit every $378$ days, purge
gap $H$ days between last training label and first test date. \textbf{Benchmarks:} Historical Simulation;
EWMA (RiskMetrics, $\lambda=0.94$) Gaussian; EVT peaks-over-threshold GPD (threshold at the window $90\%$
quantile); GARCH(1,1)-$t$ refit every $126$ days with an $O(N)$ variance filter between refits.

\bigskip
\noindent\emph{Note on citation checks:} the pre-2025 references were verified against publisher records;
the 2026 references were confirmed to exist. Before submission, verify the exact \emph{Finance Research
Letters} volume/article number and DOI for Li et al.\ (2026) and the final volume/pages for the Journal of
Forecasting item on the publisher page.

\end{document}
